\subsection{VECTOR BUNDLES WITH OPERATORS}

Let G be a Lie group, X a manifold of class \(C^{r}\) and \((g, x) \mapsto gx\) a law of left operation of class \(C^{r}\) of G on X. Let E be a vector bundle of class \(C^{r}\) , with base space X and \(\pi: E \to X\) the projection of E onto X. For all \(x \in X\) , let \(E_{x}\) be the fibre of E at x. Let \((g, u) \mapsto gu\) be a law of left operation of G on E such that \(\pi\) is compatible with the operations of G on X and on E. For all \(g \in G\) and all \(x \in X\) , the restriction to \(E_{x}\) of the mapping \(u \mapsto gu\) is a bijection \(\psi_{g,x}\) of \(E_{x}\) onto \(E_{gx}\) . We shall assume that, for all \(g \in G\) and all \(x \in X\) , \(\psi_{g,x}\) is continuous and linear and hence is an isomorphism of the normable space \(E_{x}\) onto the normable space \(E_{gx}\) .

Let \(\phi\) be the automorphism \((g,x)\mapsto (g,gx)\) of the manifold \(\mathrm{G}\times \mathrm{X}\) . Let \(p\) be the canonical projection of \(\mathrm{G}\times \mathrm{X}\) onto \(\mathbf{X}\) and \(\mathrm{E}'\) the inverse image of \(\mathrm{E}\) under \(p\) . Let \(\psi \colon \mathrm{E}'\to \mathrm{E}'\) be the mapping the sum of the \(\psi_{g,x}\colon \mathrm{E}_{(g,x)}'\to \mathrm{E}_{(g,gx)}'\) .

DEFINITION 4. If \(\psi\) is a \(\phi\) -morphism of vector bundles of class \(\mathbf{C}^r\) , \(\mathbf{E}\) is called a vector \(\mathbf{G}\) -bundle of class \(\mathbf{C}^r\) .

In other words, E is a vector G-bundle of class \(C^{r}\) if for all \((g_{0}, x_{0}) \in \mathbf{G} \times \mathbf{X}\) the following condition holds: there exists an open neighbourhood U of \((g_{0}, x_{0})\) in \(G \times X\) such that, if \(E^{\prime} \mid U\) (resp. \(E^{\prime} \mid \phi(U)\) ) is identified with a trivial vector bundle of fibre M (resp. N) by means of a vector chart, the mapping \((g, x) \mapsto \psi_{g,x}\) of U into \(\mathcal{L}(M, N)\) is of class \(C^{r}\) .

The mapping \(\psi\) is obviously bijective and it follows from the above local criterion that \(\psi^{-1}\) is a \(\phi^{-1}\) -morphism of vector bundles so that \(\psi\) is a \(\phi\) -isomorphism of vector bundles.

A trivial vector G-bundle of base X is a vector bundle \(X \times F\) (where F is a complete normable space) with the law of operation \((g, (x, f)) \mapsto (gx, f)\) of G on \(X \times F\) .

We again assume the hypotheses and notation preceding Definition 4 and further take \(\tau\) to be a vector functor of class \(C^r\) for isomorphisms (Differentiable and Analytic Manifolds, R, 7.6.6). Then \(\tau E\) is a vector bundle with base space \(X\) . For all \(x \in X\) , its fibre \((\tau E)_x\) is equal to \(\tau (E_x)\) . For all normable spaces \(N_1, N_2\) , let \(\operatorname{Isom}(N_1, N_2)\) denote the set of isomorphisms of \(N_1\) onto \(N_2\) . If \(g \in G\) , then

\[
\tau (\psi_ {g, x}) \in \operatorname{Isom} ((\tau E) _ {x}, (\tau E) _ {g x}).
\]

The \(\tau(\psi_{g,x})\) define a law of left operation \((g,u)\mapsto gu\) of G on \(\tau\) E and the canonical projection of \(\tau\) E onto X is compatible with the operations of G on X and \(\tau\) E.

PROPOSITION 15. If E is a vector G-bundle of class \(\mathbf{C}^r\) , \(\tau \mathbf{E}\) is a vector G-bundle of class \(\mathbf{C}^r\) .

Let \(g_0, x_0\) , U, M, N be as in the paragraph following Definition 4. Then the mapping \((g, x) \mapsto \tau(\psi_{g,x})\) of U into \(\mathcal{L}(\tau\mathrm{M}, \tau\mathrm{N})\) is the composition of the mapping \((g, x) \mapsto \psi_{g,x}\) of U into \(\mathcal{L}(\mathrm{M}, \mathrm{N})\) and the mapping \(f \mapsto \tau(f)\) of \(\operatorname{Isom}(\mathrm{M}, \mathrm{N})\) into \(\operatorname{Isom}(\tau\mathrm{M}, \tau\mathrm{N})\) ; these two mappings are of class \(C^r\) and hence so is their composition, whence the proposition.

PROPOSITION 16. Let G be a Lie group, X a manifold of class \(\mathbf{C}^r\) ( \(r \geqslant 2\) ) and \((g, x) \mapsto gx\) a law of left operation of class \(\mathbf{C}^r\) of G on X, whence, by transporting the structure, there is a law of left operation of G on TX. Under this law, TX is a vector G-bundle of class \(\mathbf{C}^{r-1}\) .

Let \(\mathrm{pr}_1\) (resp. \(\mathrm{pr}_2\) ) be the canonical projection of \(\mathrm{G} \times \mathrm{X}\) onto \(\mathrm{G}\) (resp. \(\mathrm{X}\) ) and let \(\mathrm{E}_1\) (resp. \(\mathrm{E}_2\) ) be the inverse image of TG (resp. TX) relative to \(\mathrm{pr}_1\) (resp. \(\mathrm{pr}_2\) ). Then the vector bundle \(\mathrm{T}(\mathrm{G} \times \mathrm{X})\) is the direct sum of \(\mathrm{E}_1\) and \(\mathrm{E}_2\) . Let \(i: \mathrm{E}_2 \to \mathrm{T}(\mathrm{G} \times \mathrm{X})\) and \(q: \mathrm{T}(\mathrm{G} \times \mathrm{X}) \to \mathrm{E}_2\) be the canonical vector bundle morphisms defined by this decomposition into a direct sum. Let \(\phi\) be the mapping \((g, x) \mapsto (g, gx)\) of \(\mathrm{G} \times \mathrm{X}\) into \(\mathrm{G} \times \mathrm{X}\) . Then the mapping denoted by \(\psi\) in Definition 4 (where we put \(\mathrm{E} = \mathrm{TX}\) ) is just \(q \circ \mathrm{T}(\phi) \circ i\) . But \(\mathrm{T}(\phi)\) is a \(\phi\) -morphism of vector bundles of class \(\mathrm{C}^{r-1}\) (Differentiable and Analytic Manifolds, R, 8.1.2).

COROLLARY. If \(\tau\) is a vector functor of class \(\mathbf{C}^r\) for isomorphisms, \(\tau(\mathrm{TX})\) is a vector G-bundle of class \(\mathbf{C}^{r-1}\) .

This follows from Propositions 15 and 16.

Remark 1. If \(\tau\) is a vector functor of class \(C^r\) for isomorphisms in finite dimension and \(E\) is of finite rank, \(\tau E\) is defined similarly and Proposition 15 remains valid; the Corollary to Proposition 16 remains valid provided \(X\) is finite-dimensional.

Examples. With the hypotheses and notation of Proposition 16, let F be a complete normable space. Then \(\mathcal{L}((\mathrm{TX})^p; \mathrm{F})\) is a vector G-bundle of class \(\mathrm{C}^{r-1}\) ; so is \(\operatorname{Alt}^p(\mathrm{TX}; \mathrm{F})\) if K is of characteristic zero or X is finite-dimensional (cf.~Differentiable and Analytic Manifolds, R, 7.7, 7.8). If X is finite-dimensional, \(\bigotimes^p(\mathrm{TX}) \otimes \bigotimes^q(\mathrm{TX})^*\) is a vector G-bundle of class \(\mathrm{C}^{r-1}\) .

PROPOSITION 17. Let G be a Lie group, X a left Lie homogeneous space of G, \(x_0\) a point of X, \(G_0\) the stabilizer \(x_0\) in G, E and E' left vector G-bundles of class \(C^{r}\) and base space X, \(E_0\) (resp. \(E_0'\) ) the fibre of E (resp. E') at \(x_0\) and f an element of \(\mathcal{L}(E_0, E_0')\) such that \(f(gu) = g(f(u))\) for all \(u \in E_0\) and \(g \in G_0\) . Then there exists one and only one morphism of E into E' compatible with the operations of G and extending f. The uniqueness of this morphism is obvious. We prove its existence. Let \(g\) , \(g'\) elements of G and \(u \in E_0\) be such that \(gu = g'u\) . Then \(g'^{-1}g \in G_0\) and \(g'^{-1}gu = u\) and hence \(g'^{-1}gf(u) = f(u)\) , that is \(gf(u) = g'f(u)\) . Hence a mapping \(\phi\) is defined of E into \(E'\) by writing \(\phi(gu) = gf(u)\) . Clearly this mapping extends \(f\) and it is compatible with the operations of G. We show that \(\phi\) is a vector bundle morphism of class \(C^r\) . Let \(x_1 \in X\) . There exists an open neighbourhood V of \(x_1\) in X and a submanifold W of G such that the mapping \(g \mapsto gx_0\) is an isomorphism \(\theta\) of class \(C^r\) of W onto V. By shrinking V and W it can be assumed that:

\begin{enumerate}
\def\labelenumi{(\arabic{enumi})}
\item
  E \textbar{} V (resp. E' \textbar{} V) is identified with a trivial vector bundle of fibre M (resp. M');
\item
  if \(\psi_{g}\) (resp. \(\psi_{g}^{\prime}\) ) denotes the mapping \(u\mapsto gu\) of \(\mathrm{E}_0\) (resp. \(\mathrm{E}_0^{\prime}\) ) into \(\mathrm{E}_{gx_0}\) (resp. \(\mathrm{E}_{gx_0}\) ), then the mappings \(g\mapsto \psi_g\) and \(g\mapsto \psi_g^{-1}\) (resp. \(g\mapsto \psi_g^{\prime}\) and \(g\mapsto \psi_g^{\prime -1}\) ) of W into \(\mathcal{L}(\mathrm{E}_0,\mathrm{M})\) and \(\mathcal{L}(\mathrm{M},\mathrm{E}_0)\) (resp. \(\mathcal{L}(\mathrm{E}_0',\mathrm{M}')\) and \(\mathcal{L}(\mathrm{M}',\mathrm{E}_0'))\) are of class \(C^r\) .
\end{enumerate}

For \(x \in V\) , let \(\phi_x: M \to M'\) be the restriction of \(\phi\) to \(E_x = M\) . Then \(\phi_x\) is obtained by composing the following mappings:

\begin{enumerate}
\def\labelenumi{(\arabic{enumi})}
\item
  the mapping \((\psi_{\theta^{-1}(x)})^{-1}\) of M into \(E_{0}\) ;
\item
  the mapping f of \(E_{0}\) into \(E_{0}'\) ;
\item
  the mapping \(\psi_{\theta^{-1}(x)}^{\prime}\) of \(\mathbf{E}_0'\) into \(\mathbf{M}'\) .
\end{enumerate}

Hence we see that the mapping \(x \mapsto \phi_x\) of V into \(\mathcal{L}(\mathbf{M}, \mathbf{M}')\) is of class \(\mathbf{C}^r\) .

COROLLARY 1. Let \(\mathrm{E}_0^{\mathbb{G}_0}\) be the set of elements of \(\mathrm{E}_0\) which are invariant under \(\mathbf{G}_0\) . For all \(u \in \mathrm{E}_0^{\mathbb{G}_0}\) , let \(\sigma_u\) be the mapping of X into E defined by \(\sigma_u(gx_0) = gu\) for all \(g \in G\) .

\begin{enumerate}
\def\labelenumi{(\roman{enumi})}
\item
  The G-invariant sections† of E are of class \(\mathbf{C}^r\) .
\item
  \(u \mapsto \sigma_u\) is a bijection of \(\mathrm{E}_0^{\mathrm{G}_0}\) onto the set of G-invariant sections of E.
\end{enumerate}

Assertion (ii) is obvious. To prove (i) it is sufficient to prove that each section \(\sigma_{u}\) is of class \(C^r\) . Let E' be the trivial G-bundle of base X and fibre \(\mathrm{E}_0^{\mathrm{G}_0}\) . Let \(f\) be the canonical injection of \(\mathrm{E}_0^{\mathrm{G}_0}\) into \(\mathrm{E}_0\) . By Proposition 17 there exists a morphism \(\phi\) of E' into E compatible with the operations of G and extending \(f\) . If \(u \in \mathrm{E}_0^{\mathrm{G}_0}\) and \(g \in \mathrm{G}\) , then

\[
\sigma_ {u} (g x _ {0}) = g u = g f (u) = \phi (g u) = \phi ((u, g x _ {0}))
\]

and hence \(\sigma_u(x) = \phi((u, x))\) for all \(x \in \mathbf{X}\) , which proves our assertion.

*For example, let G be a finite-dimensional real Lie group, \(G_{0}\) a compact Lie subgroup of G and X the homogeneous space \(G/G_{0}\) . Let \(x_{0}\) denote the canonical image of e in X. There exists a positive definite symmetric bilinear form on \(T_{x_{0}}(X)\) invariant under \(G_{0}\) (Integration, Chapter VII, § 3, Proposition 1). Applying the above to \((\mathrm{TX})^{*}\otimes(\mathrm{TX})^{*}\) , we see that there exists on X an analytic Riemannian metric invariant under \(G_{0}\) .

COROLLARY 2. Suppose that \(G_{0}\) operates trivially on \(E_{0}\) . Let \(E'\) be the trivial G-bundle with base space X and fibre \(E_{0}\) . There exists one and only one isomorphism of E onto \(E'\) compatible with the operations of G and extending \(Id_{E_{0}}\) .

This follows immediately from Proposition 17.

Remark 2. In this no., the laws of left operation can be replaced throughout by laws of right operation.

